\chapter{Bruxism}

\section*{Definition}
Bruxism is the repetitive masticatory muscle activity characterized by clenching or grinding of the teeth, occurring during sleep (sleep bruxism) or wakefulness (awake bruxism), with or without audible grinding sounds.

\section*{When to See a Doctor}

\begin{itemize}
 \item Bed partner reports loud grinding sounds, especially if associated with choking/gasping (suggests OSA)
 \item Waking with jaw pain, temporal headache, or tooth sensitivity
 \item Visible tooth wear (flattening, chipping, fractures) noticed by dentist or patient
 \item Difficulty opening the mouth or locking (may indicate temporomandibular joint involvement)
\end{itemize}

\section*{How It Develops}
Sleep bruxism is associated with brief sleep arousals and activity of the jaw muscles, while awake bruxism involves repeated or sustained jaw activity during wakefulness. Stress, concentration, sleep disruption, medicines, and other factors may contribute; the exact cause varies.

\section*{Causes}
\begin{itemize}
 \item \textbf{Psychosocial:} Acute or chronic stress, anxiety disorders, type~A personality, suppressed anger or aggression
 \item \textbf{Sleep-related:} Obstructive sleep apnea (7--8x increased odds), REM sleep behavior disorder, periodic limb movements, insomnia
 \item \textbf{Pharmacologic:} Selective serotonin reuptake inhibitors (SSRIs) and serotonin--norepinephrine reuptake inhibitors (SNRIs), amphetamines, cocaine, MDMA, bupropion, antipsychotics (especially haloperidol), caffeine, nicotine, alcohol withdrawal
 \item \textbf{Neurologic:} Dystonia (oromandibular), Huntington disease, Parkinson disease, traumatic brain injury, cerebellar degeneration
 \item \textbf{Dental/Malocclusion:} Controversial as primary cause; may exacerbate existing bruxism
\end{itemize}

\section*{Related Symptoms}
\begin{itemize}
 \item Myofascial pain (masseter, temporalis, medial pterygoid), morning headache (typically temporal), tooth hypersensitivity, TMJ clicking or locking, hypertrophied masseter muscles, tongue scalloping, cheek biting, cervical muscle tension, tinnitus
\end{itemize}
\textbf{Important:} Sleep bruxism is a risk indicator for obstructive sleep apnea; screening for OSA is essential before initiating any treatment that may mask this underlying condition.

\section*{Medical Evaluation}
\begin{itemize}
 \item Clinical history: onset, frequency, duration; partner report of grinding sounds; morning symptoms; medication and substance use review
 \item Dental examination: assess tooth wear (incisal/occlusal facets), enamel fractures, abfraction lesions, gingival recession, hypertrophied masseters
 \item Sleep testing is considered when obstructive sleep apnea is suspected, such as with loud snoring, witnessed pauses, gasping, or excessive daytime sleepiness. Instrumental testing for bruxism is usually reserved for selected cases.
 \item Clinical questionnaires: Bruxism Episode Index, Oral Behavior Checklist, TMD screening
\end{itemize}

\section*{Treatment Options}
\begin{itemize}
 \item \textbf{Occlusal splint therapy:} Custom-fitted hard acrylic or soft mouthguard worn at night to protect teeth and reduce grinding force; does not stop bruxism but mitigates dental damage
 \item \textbf{Behavioral therapy:} CBT for stress reduction, biofeedback, hypnosis, sleep hygiene optimization
 \item \textbf{Medicines and injections:} There is no compelling evidence for routine drug treatment of sleep bruxism. Medicines or botulinum toxin are specialist, individualized, and usually off-label options because benefits and risks remain uncertain.
 \item \textbf{Treat underlying sleep disorders:} Assessment and treatment of obstructive sleep apnea may improve symptoms when both conditions are present
 \item \textbf{Dental restoration:} Repair fractured or worn teeth when needed; irreversible bite adjustment should not be done solely to treat bruxism
\end{itemize}

\section*{Self-Care and Home Management}
\begin{itemize}
 \item Apply warm moist heat to the masseter and temporalis muscles for 10--15 minutes before sleep
 \item Practice jaw relaxation exercises: let the jaw hang with lips together, teeth apart (rest position)
 \item Avoid chewing gum, ice, pens, or other non-food objects during the day
 \item Reduce evening caffeine (coffee, tea, chocolate) and alcohol intake
 \item Establish a calming bedtime routine: meditation, deep breathing, progressive muscle relaxation
 \item Sleep with a properly fitted occlusal splint (night guard) if recommended by a dentist
\end{itemize}

\section*{References}
\begin{thebibliography}{99}
\bibitem{bruxism:ref1} Lobbezoo F, Ahlberg J, Glaros AG, et al. ``Bruxism Defined and Graded: An International Consensus.'' J Oral Rehabil. 2013;40(1):2--4.
\bibitem{bruxism:ref2} Manfredini D, Restrepo C, Diaz-Serrano K, et al. ``Prevalence of Bruxism in Children and Adults: A Systematic Review.'' J Oral Rehabil. 2013;40(9):685--702.
\bibitem{bruxism:ref3} Murali RV, Rangarajan P, Mounissamy A. ``Bruxism: Conceptual Discussion and Review.'' J Pharm Bioallied Sci. 2015;7(Suppl 1):S265--S267.
\bibitem{bruxism:ref4} Ohayon MM, Li KK, Guilleminault C. ``Risk Factors for Sleep Bruxism in the General Population.'' Chest. 2001;119(1):53--61.
\bibitem{bruxism:ref5} Carra MC, Huynh N, Lavigne G. ``Sleep Bruxism: A Comprehensive Overview for the Dental Clinician.'' J Can Dent Assoc. 2012;78:c64.
\bibitem{bruxism:ref6} Shetty S, Pitti V, Babu CLS, et al. ``Bruxism: A Literature Review.'' J Indian Prosthodont Soc. 2010;10(3):141--148.
\end{thebibliography}
